Temporal information in an autonomous quantum system need not be accompanied by global time-translation asymmetry: a clock and signal can carry a relative temporal mode while their joint state remains inside a fixed total-energy sector. What physical resource then limits the information in that mode? We first show that local tangent Fisher information alone has no universal high-frequency energy bound: at fixed mean energy it can remain finite while the affine neighborhood over which its tangent is physical collapses. This exposes two complementary regimes. At nonzero physical tangent radius, arbitrary finite-copy collective measurements obey a two-sided spectral-survival bound. At a rank-changing boundary the resource changes order: positive second-order endpoint synthesis is required, yielding sharp common-record action bounds with coefficients $\hbar\nu/4$ for bilateral exchange and $\hbar\nu/2$ for one-sided synthesis. The same clean boundary action also bounds the SLD-QFI trace with sharp coefficient $\hbar\nu/4$. By contrast, the additional rank-changing Bures contribution contains independent zero-eigenvalue Hessian curvature; a spectator construction shows that it cannot be assigned to the selected first-order temporal mode without separately pricing that curvature. The laws extend to coherent baselines and multiple exchange frequencies, and fixed-total-energy constructions saturate both the single-gap constants and the complete frequency-weighted sum. Thus relational temporal information remains resource constrained even when global asymmetry and signed energy curvature vanish.